\chapter*{Abstract}\label{sec:Abstract}
 \addcontentsline{toc}{chapter}{Abstract}

Electric racing motorcycles require accurate, real-time knowledge of the battery's 
State of Charge (SOC) and State of Health (SOH) to guarantee safe and reliable 
operation under the highly dynamic current profiles typical of a race. This work 
develops and validates a model-based estimation framework for the simultaneous 
estimation of SOC and SOH of the WeLion SHP270-30 Nickel-Manganese-Cobalt (NMC) 
lithium-ion cell (30~Ah nominal capacity) used in the eRacing electric 
motorcycle developed at ETSEIB.

The battery was represented using a second-order Thévenin equivalent circuit model 
(ECM), consisting of a voltage source, a series resistance, and 
two parallel resistor--capacitor branches. The cell was first characterised 
experimentally through discharge, charge, and pulse-discharge tests, from which the 
actual capacity and the Open Circuit Voltage-SOC relationship were obtained. The 
equivalent-circuit parameters were subsequently identified using an algebraic method, 
applied to a two-level step-current test repeated across different 
SOC operating points. Finally, two Extended Kalman Filters (EKF) were 
implemented in MATLAB, a SOC filter, and a dual EKF that simultaneously estimates 
SOC and the SOH parameters, the series resistance $R_0$ and the 
capacity $Q$.

Both filters were validated using experimental data from three current profiles: a 
charge cycle, a cyclic pulse discharge, and a random step-current discharge simulating 
race conditions. The SOC filter achieved root-mean-square errors (RMSE) between 
$0.69\%$ and $1.14\%$ relative to a coulomb-counting reference, remaining within the 
estimated $3\sigma$ uncertainty bounds throughout all tests. The dual EKF achieved 
comparable SOC accuracy (RMSE between $0.94\%$ and $1.39\%$) while simultaneously 
obtaining the capacity and resistance estimates to values consistent with the 
characterised cell.

The results confirm that a second-order Thévenin ECM combined with a dual EKF 
provides a computationally efficient and sufficiently accurate framework for 
real-time SOC and SOH estimation under the dynamic current profiles found in 
motorsport applications. The algebraic parameter identification method proved 
effective across most of the operating range, although reduced reliability was 
observed at the extremes of the state of charge, consistent with the known numerical 
sensitivity of the method and the limitations of the equivalent-circuit model in 
capturing the underlying electrochemical behaviour in these regions. Future work 
should extend the framework to account for temperature dependence and validate its 
performance under real race data.

\vspace{1em}
\noindent\textbf{Keywords:} State of Charge, State of Health, Extended Kalman Filter, 
Equivalent-Circuit Model, lithium-ion battery, battery characterisation.


\chapter*{Resum}
\addcontentsline{toc}{chapter}{Resum}
Les motocicletes elèctriques de competició requereixen un coneixement precís i en 
temps real de l'Estat de Càrrega (SOC) i l'Estat de Salut (SOH) de la bateria per 
garantir un funcionament segur i fiable sota els perfils de corrent altament dinàmics 
propis d'una cursa. Aquest treball desenvolupa i valida un marc d'estimació basat en un 
model per a l'estimació simultània del SOC i el SOH de la cel·la de liti-ió WeLion 
SHP270--30, de química Níquel-Manganès-Cobalt (NMC) i capacitat nominal de 30~Ah, 
utilitzada a la motocicleta elèctrica eRacing desenvolupada a l'ETSEIB.

La bateria es va representar mitjançant un model de circuit equivalent (ECM) de 
Thévenin de segon ordre, format per una font de tensió, una resistència en sèrie i 
dues branques paral·leles de resistència-condensador. La cel·la es va caracteritzar 
experimentalment mitjançant assajos de descàrrega contínua, càrrega i descàrrega per polsos, 
dels quals es van obtenir la capacitat real i la relació OCV--SOC. Posteriorment, els 
paràmetres del circuit equivalent es van identificar mitjançant un mètode algebraic, 
aplicat a un assaig d'esglaó de corrent de dos nivells, repetit en diferents punts 
d'operació del SOC. Finalment, es van implementar dos Filtres de Kalman Estesos (EKF) 
en MATLAB, un filtre de SOC, i un EKF dual que estima simultàniament el SOC i els 
paràmetres del SOH, la resistència en sèrie $R_0$ i la capacitat $Q$.

Ambdós filtres es van validar amb dades experimentals de tres perfils de corrent: un 
cicle de càrrega, una descàrrega cíclica per polsos, i una descàrrega d'esglaons de 
corrent aleatoris que simulen les condicions d'una cursa. El filtre de SOC va assolir 
errors quadràtics mitjans (RMSE) d'entre el $0,69\%$ i l'$1,14\%$ respecte a una 
referència de comptatge de coulombs, mantenint-se dins dels límits d'incertesa de 
$3\sigma$ estimats en tots els assajos. L'EKF dual va assolir una precisió del SOC 
comparable (RMSE d'entre el $0,94\%$ i l'$1,39\%$) mentre convergia simultàniament 
les estimacions de capacitat i resistència a valors consistents amb la cel·la 
caracteritzada.

Els resultats confirmen que un ECM de Thévenin de segon ordre combinat amb un EKF dual 
proporciona un marc computacionalment eficient i prou precís per l'estimació en 
temps real del SOC i el SOH sota els perfils de corrent dinàmics propis de les 
aplicacions d'automobilisme esportiu. El mètode algebraic d'identificació de 
paràmetres es va mostrar efectiu en la major part del rang d'operació, tot i que es va 
observar una fiabilitat reduïda als extrems de l'estat de càrrega, coherent amb la 
sensibilitat numèrica coneguda del mètode i amb les limitacions del model de circuit 
equivalent per capturar el comportament electroquímic subjacent en aquestes regions. 
Els treballs futurs haurien d'ampliar el marc per tenir en compte la dependència de la 
temperatura i validar-ne el rendiment amb dades reals de telemetria d'una cursa.

\vspace{1em}
\noindent\textbf{Paraules clau:} Estat de Càrrega, Estat de Salut, Filtre de Kalman 
Estès, Model de Circuit Equivalent, bateria de liti-ió, caracterització de bateries.


\chapter*{Resumen}
\addcontentsline{toc}{chapter}{Resumen}
Las motocicletas eléctricas de competencia requieren un conocimiento del Estado de Carga (SOC) preciso y en 
tiempo real, así como el Estado de Salud (SOH) de la batería para poder 
garantizar un funcionamiento seguro y confiable bajo perfiles dinámicos de corriente 
correspondientes a los de una carrera. Este trabajo desarrolla y valida una metodología
para la estimación simultánea del SOC y el SOH de la celda de litio-ion WeLion SHP270--30, de química 
Níquel-Manganeso-Cobalto (NMC), con una capacidad nominal de $30 Ah$, 
la cual es la correspondiente a la utilizada en la motocicleta eléctrica eRacing desarrollada 
en ETSEIB.

La batería se representó mediante un modelo de circuito equivalente (ECM) de Thévenin 
de segundo orden, conformado por una fuente de voltaje, una resistencia en serie y dos 
ramas con una resistencia y un capacitor en paralelo. La celda se caracterizó experimentalmente 
mediante pruebas de descarga, carga y descarga por pulsos, de las cuales se obtuvieron 
la capacidad real y la relación OCV--SOC. Posteriormente, los parámetros del circuito 
equivalente se identificaron mediante un método algebraico, aplicado a una prueba 
conformada por dos escalones de corriente y repetida en distintos puntos de operación del SOC. 
Finalmente, se implementaron dos filtros de Kalman Extendidos (EKF) en MATLAB, un 
filtro de SOC, y un dual EKF que estima simultáneamente el SOC y los parámetros del 
SOH, siendo estos parámetros la resistencia en serie $R_0$ y la capacidad $Q$.

Ambos filtros se validaron con datos experimentales obtenidos de tres perfiles de corriente: un 
ciclo de carga, uno de descarga cíclica por pulsos, y uno de descarga de escalones de 
corriente aleatorios que simulan las condiciones de una carrera. El filtro de SOC 
alcanzó errores (RMSE) de entre $0.69\%$ y $1.14\%$ respecto a la 
referencia de Coulomb Counting, manteniéndose dentro de los límites de incertidumbre 
de $3\sigma$ estimados en todas las pruebas. El EKF dual alcanzó una precisión del SOC 
comparable (RMSE de entre $0.94\%$ y $1.39\%$) mientras obtenía al mismo tiempo las 
estimaciones de capacidad y resistencia a valores consistentes con la celda 
caracterizada.

Los resultados confirman que un \acs{ECM} de Thévenin de segundo orden combinado con un dual \acs{EKF} 
proporciona un marco computacionalmente eficiente y suficientemente preciso para 
la estimación en tiempo real del SOC y el SOH bajo perfiles de corriente dinámicos aproximadamente
caracteristicos de las aplicaciones del automovilismo deportivo. El método algebraico de 
identificación de parámetros demostró ser efectivo en la mayor parte del rango de 
operación, aunque se observó una confiabilidad reducida en los extremos del estado de 
carga, congruente con la sensibilidad numérica conocida del método y con las 
limitaciones del modelo de circuito equivalente para capturar el comportamiento 
electroquímico en estas regiones. Futuros trabajos deberían de ampliar el 
marco para considerar la dependencia con la temperatura y validar su desempeño con 
datos reales de telemetría de carrera.

\vspace{1em}
\noindent\textbf{Palabras clave:} Estado de Carga, Estado de Salud, Filtro de Kalman 
Extendido, Modelo de Circuito Equivalente, batería de litio-ion, caracterización de 
baterías.
